% authority: science_draft
% claim_status: draft/control
\documentclass[11pt]{article}
\usepackage[margin=1in]{geometry}
\usepackage{amsmath,amssymb,amsthm,booktabs,longtable,array,enumitem,hyperref}
\hypersetup{hidelinks}
\setlist{nosep}
\newtheorem{definition}{Definition}
\newtheorem{proposition}{Proposition}
\newtheorem{lemma}{Lemma}
\newcommand{\A}{\mathcal A}
\newcommand{\Q}{\mathcal Q_{\mathrm{B1}}}
\newcommand{\Qreg}{\mathcal Q_{\mathrm{B1,reg}}}
\newcommand{\E}{E_{\mathrm{B1}}}
\newcommand{\R}{\mathcal R_{\mathrm{B1}}}
\newcommand{\sigm}{\operatorname{logistic}}

\title{V22 P3-T01 Local Multi-Field Source State v1\\
\large Exact local syntax for the sealed B1 research candidate}
\author{Candidate Constructor with Ontology-Formalizer perspective}
\date{2026-08-09}

\begin{document}
\sloppy
\maketitle

\begin{center}
\fbox{\parbox{0.92\textwidth}{
\textbf{Status and authority.} This is a proposal-only draft/control
construction. It activates the sealed B1 packet only as the sole research
candidate eligible for P3-T02. It does not adopt the bundle or a source law,
modify canonical ontology, interpret the source arena as physical spacetime,
construct a causal cone or effective metric, establish realistic matter or
universal coupling, pass Gate B, or complete a GR derivation.
}}
\end{center}

\section{Exact object and source/target boundary}

Let $\A$ be the already-assumed smooth four-dimensional \emph{source} arena.
Its smoothness, dimension, topology, open sets, and charts are primitive debt.
They are not derived, and no source chart is a target atlas, physical position
system, clock reading, or distance assignment.  P3-T01 constructs the sealed
B1 candidate
\[
   \E=\A\times\mathbb R^6,\qquad \pi:\E\to\A,
\]
as a proposal-only trivial rank-six real vector bundle.  No metric, coframe,
connection, orientation, volume form, causal order, or preferred chart is part
of $\E$.

\begin{definition}[B1 local-state sheaf]
For each source-open $U\subseteq\A$, define
\[
  \Q(U)=\Gamma_{C^2}(U,\E|_U)=C^2(U,\mathbb R^6),
  \qquad q=(q^1,\ldots,q^6).
\]
For $V\subseteq U$, the restriction is
$\operatorname{res}_{UV}(q)=q|_V$ componentwise.  The six components are
source scalars.  The selected structure is a sheaf because P3-T01 needs exact
restriction, overlap agreement, and unique local-to-global gluing.  A cosheaf
or chain complex would not express those obligations without additional
structure.
\end{definition}

The components $q^1,q^2,q^6$ share a declared source unit so the ratios below
are dimensionless.  The other component units remain declared source units;
none is a target length, duration, mass, or charge.  $C^2$ regularity supplies
the first jets required by the sealed B1 identity and one additional derivative
for stable local variation tests.  It is not evidence for dynamics.

\section{Restriction, gluing, and regular operational locus}

The restriction laws are
\[
 \operatorname{res}_{UU}=\mathrm{id},\qquad
 \operatorname{res}_{VW}\circ\operatorname{res}_{UV}
 =\operatorname{res}_{UW}\quad(W\subseteq V\subseteq U).
\]
For an open cover $U=\bigcup_i U_i$, compatible sections satisfy
\[
 q_i|_{U_i\cap U_j}=q_j|_{U_i\cap U_j}
 \quad\text{componentwise for all }i,j.
\]
Ordinary smooth-function gluing produces one and only one $q\in\Q(U)$ with
$q|_{U_i}=q_i$.  The fiber transition is the identity in the fixed
$\mathbb R^6$ fiber.  Source-chart representatives transform as six scalars;
they are not glued with target coordinates.

On the denominator locus $q^6\ne0$, set
\[
  \rho_1=\frac{q^1}{q^6},\qquad
  \rho_2=\frac{q^2}{q^6},\qquad
  p_i=\sigm(\rho_i)=\frac{1}{1+e^{-\rho_i}}.
\]
The regular operational subdomain is
\[
 \Qreg(U)=\{q\in\Q(U):q^6(x)\ne0,
   (d\rho_1\wedge d\rho_2)(x)\ne0\ \text{for every }x\in U\}.
\]
This pointwise condition is stable under restriction and under compatible
gluing by a regular cover.  Sections can still exist where $q^6=0$ or
$d\rho_1\wedge d\rho_2=0$, but the two-channel interface is respectively
undefined or rank deficient there.  No silent regularization is allowed.

\begin{proposition}[Sheaf and observable naturality]
The assignments $\Q$, $\rho_i$, $p_i$, and $d\rho_i$ commute with source-open
restriction on their declared domains.  Compatible regular local sections
glue uniquely to a regular section.
\end{proposition}
\begin{proof}
The sheaf axioms are the componentwise $C^2$ function axioms.  Division by a
nowhere-zero scalar, smooth logistic composition, and exterior differentiation
are local operations, so each commutes with restriction.  Compatibility makes
the local ratios and their exterior derivatives equal on overlaps.  Their
pointwise nonvanishing conditions therefore pass to the unique glued section.
\end{proof}

Three fail-closed controls distinguish the claims.  A component mismatch on
an overlap admits no glued base section.  Compatible base sections whose
glued $q^6$ vanishes do glue in $\Q$ but not in $\Qreg$.  A rule that
uses target-coordinate Jacobians, a coframe, or a metric to glue the six
scalars is not this candidate.

\section{Variations and boundary data}

For a declared source pair $(U,B)$, let $i_B:B\hookrightarrow\overline U$ be
an embedded closed source hypersurface or patch-boundary stratum.  The
available boundary data are the value trace $i_B^*q$ and tangential first-jet
trace $i_B^*(dq)$.  Neither a unit normal, normal derivative, induced metric,
boundary measure, nor causal-boundary meaning is defined.

Boundary-preserving variations are compactly supported $C^2$ vertical
sections $\delta q$ satisfying
\[
 i_B^*(\delta q)=0,\qquad i_B^*(d\delta q)=0.
\]
For each relatively compact $V\Subset U$ and $q\in\Qreg(U)$,
sufficiently small $C^1$ variations keep the two nonvanishing conditions on
$V$.  This is an ordinary open-condition statement in local source charts;
it assumes no metric norm on $\A$.

\begin{proposition}[Two independent local response directions]
At every $x_0\in U$ and $q\in\Qreg(U)$, the vertical differential of
\[
 \R(q)=(p_1,p_2)
\]
has rank at least two on admissible compactly supported variations.
\end{proposition}
\begin{proof}
Choose a $C^2$ bump function $\chi$ supported away from $B$ with
$\chi(x_0)=1$.  Let $\delta_1q^1=\chi q^6$ with all other components zero,
and let $\delta_2q^2=\chi q^6$ with all other components zero.  At $x_0$,
\[
 \delta_1(\rho_1,\rho_2)=(1,0),\qquad
 \delta_2(\rho_1,\rho_2)=(0,1).
\]
Because $\sigm'(z)=\sigm(z)(1-\sigm(z))>0$ for finite $z$, the corresponding
response vectors are linearly independent.  Thus the local response rank is
two.  This clears only the P2-T01 necessary information-capacity screen; it is
not a metric, cone, scale, or Gate-B sufficiency result.
\end{proof}

\section{Source representation equivalence}

The source-chart transition pseudogroup acts passively.  If $\phi$ changes a
source representation, a local scalar representative acts by pullback
$q\mapsto\phi^*q$.  The ratios and token parameters are scalars, while
$d\rho_i$ and $d\rho_1\wedge d\rho_2$ transform by the corresponding tensor
pullback.  Restriction commutes with this action.  Two chart records are
equivalent exactly when they represent the same abstract section and natural
tensors on the common source domain.

No internal $GL(6)$ gauge is declared: it would generally change the selected
ratio channels.  No active dynamical source symmetry is inferred before
P3-T02 supplies a law.  Passive source representation equivalence is not a
target diffeomorphism claim or an assertion of physical gauge redundancy.

\section{Source-native observables and protected P7 interface}

The exact catalog consists of the dimensionless source scalars $\rho_1$ and
$\rho_2$, their ordered token-parameter pair $(p_1,p_2)$, and the source
two-form $d\rho_1\wedge d\rho_2$ used only as a regularity/rank witness.  None
is a physical position, clock, rod, signal, detector, causal cone, metric
volume, shear, curvature, or polarization observable.

The P7 package is consumed without modification as a protected conditional
input.  For a finite pointed sampling object $S=(J,a)$ with finite label set
$J$ and source-anchor map $a:J\to U$, define two conditional Bernoulli slots
by
\[
 \Pr(Y_{ji}=1\mid q,S)=p_i(q)(a(j)),\qquad i=1,2.
\]
The labels and anchors are source-side bookkeeping, not target detector
locations.  Restriction retains only anchors in the smaller source domain.
A refinement $r:J'\twoheadrightarrow J$ preserves declared persistent anchors
and copies their parameters exactly; new labels evaluate the same $p_i$.
Marginalization/reindexing onto persistent labels gives the parent kernel.

This interface contains no metric interpolation, nearest-neighbor rule,
volume average, convergence topology, error bound, or evolution law.  P3-T03
must construct those objects if P3-T02 first supplies viable dynamics.  Finite
P7 identities therefore remain finite protected postulates, not realistic
continuum matter or target coupling.

\section{Assumption cost and no-import audit}

The construction records nine non-scalar assumption-cost entries: the
existing source-arena debt; the rank-six bundle; the $C^2$ sheaf; compatible
units for the ratio channels; the regular-locus restriction; passive chart
naturality; source-boundary traces; two conditional logistic token bridges;
and the pointed-sampling/refinement interface.  Exactly one is the inherited
arena debt.  Every new object remains proposal-only, and the bundle remains a
new-ontology-primitive candidate rather than an adopted primitive.  The
metric-equivalent input count, target-atlas input count, and GR-target-fit
count are all zero.

All formulas are coordinate-free combinations of source scalars, restriction,
pullback, and exterior differentiation.  Local coordinates may display these
objects but supply no physical target meaning.  The construction uses no
target atlas, target metric, quadratic form, coframe, connection, volume form,
time orientation, causal cone, clock, rod, target detector, target matter
action, or benchmark prediction.

\section{Distance-to-GR and decisive result}

\begin{longtable}{p{0.28\textwidth}p{0.60\textwidth}}
\toprule
Burden & P3-T01 status \\
\midrule
Source ontology primitives & B1 bundle/sheaf is a proposal-only candidate; canonical ontology unchanged \\
Source equivalence $\mathrm{EqSrc}$ & passive chart naturality is exact locally; no adopted physical equivalence law \\
$\mathrm{RetainH}$ and $\mathrm{GenH}$ & unchanged and not used \\
$\mathrm{ObsLoc}_{\rm lc}$ & two source-token parameters are locally typed; no physical localization witness \\
$\mathrm{Resp}_{\rm lc}$ & two independent vertical response directions constructed; necessary only \\
$M_{\rm src}$ & unchanged; $\A$ is primitive source arena, not reconstructed spacetime \\
$g_{\rm eff}$ & active milestone remains blocked by absent source dynamics and reduced principal geometry \\
Matter coupling & P7 remains conditional finite input, not target coupling \\
Einstein equations & blocked; no action or target field equation \\
Finite-variation robustness & local small-variation rank stability only; dynamics/global robustness open \\
Benchmark promotion and Gate Chair & blocked and human-gated \\
Route freeze & scalar and graph routes remain frozen; exact B1 packet advances to P3-T02 only \\
\bottomrule
\end{longtable}

The decisive result is a constructed candidate: the exact sealed B1
bundle/sheaf, its domains, restrictions, gluing, regular locus, admissible
variations, boundary traces, source-representation equivalence, observables,
rank-two witness, P7 pointed-sampling interface, and assumption ledger.  The
construction discharges P3-T01's syntax burden and activates this packet only
for research execution.  It does not discharge the metric law.  P3-T02 must
construct or precisely fail a source-native dynamical law; its mathematical
refuter perspective must reject any hidden measure, derivative contraction,
connection, coefficient fit, or collapse to one response amplitude.

\end{document}
